All ablations change one setting of the SimCLR-style pretraining and keep the probe, splits and seeds fixed; they are run on the same five seeds. The full method (both augmentation families, $\tau=0.5$, 60 epochs) is the \texttt{SimCLR-style probe} row of Table~\ref{tab:main}.

\begin{table}[t]\centering\caption{Augmentation ablation: test accuracy of SimCLR with one augmentation family only (5 seeds).}\label{tab:abl}
\resizebox{0.85\columnwidth}{!}{\begin{tabular}{llcc}
\toprule
Method & Task & test\_acc $\uparrow$ & $n$ \\
\midrule
SimCLR geom-only & n50 & \textbf{0.927 $\pm$ 0.018} & 5 \\
SimCLR photo-only & n50 & \underline{0.918 $\pm$ 0.012} & 5 \\
\midrule
SimCLR geom-only & n10 & \textbf{0.725 $\pm$ 0.116} & 5 \\
SimCLR photo-only & n10 & \underline{0.702 $\pm$ 0.063} & 5 \\
\midrule
SimCLR geom-only & n200 & \textbf{0.975 $\pm$ 0.010} & 5 \\
SimCLR photo-only & n200 & \underline{0.965 $\pm$ 0.016} & 5 \\
\bottomrule
\end{tabular}
}\end{table}

\begin{table}[t]\centering\caption{Supervised scratch with the full SimCLR augmentation applied to the labeled examples (5 seeds).}\label{tab:supaug}
\resizebox{0.85\columnwidth}{!}{\begin{tabular}{llcc}
\toprule
Method & Task & test\_acc $\uparrow$ & $n$ \\
\midrule
Supervised + aug & n50 & \textbf{0.923 $\pm$ 0.033} & 5 \\
\midrule
Supervised + aug & n10 & \textbf{0.713 $\pm$ 0.099} & 5 \\
\midrule
Supervised + aug & n200 & \textbf{0.969 $\pm$ 0.014} & 5 \\
\bottomrule
\end{tabular}
}\end{table}

\textbf{H4 (augmentations matter): partly supported.} At 50 labels the full augmentation reaches \rhval{main/simclr-style-probe/n50/test_acc/mean:3} (Table~\ref{tab:main}), geometric-only \rhval{abl_aug/simclr-geom-only/n50/test_acc/mean:3} and photometric-only \rhval{abl_aug/simclr-photo-only/n50/test_acc/mean:3} (Table~\ref{tab:abl}); the drops are \rhval{cmp/cmp_aug/simclr-geom-only/n50/test_acc/delta:3} (geometric-only; paired $p=\rhval{cmp/cmp_aug/simclr-geom-only/n50/test_acc/paired_p:2}$, not resolved with five seeds) and \rhval{cmp/cmp_aug/simclr-photo-only/n50/test_acc/delta:3} (photometric-only; $p=\rhval{cmp/cmp_aug/simclr-photo-only/n50/test_acc/paired_p:3}$) (Table~\ref{tab:extra}). Only the photometric drop is clearly resolved, so H4 is supported for photometric-only and inconclusive for geometric-only. At 10 labels the ordering is the same, but the seed std of geometric-only is \rhval{abl_aug/simclr-geom-only/n10/test_acc/std:3}, so that comparison is inconclusive. Geometric-only is not distinguishable from photometric-only.

\textbf{Is the gain just augmentation?} Table~\ref{tab:supaug} gives scratch training with the same augmentations. It is not distinguishable from plain scratch at any budget (10 labels: \rhval{abl_supaug/supervised-+-aug/n10/test_acc/mean:3} vs \rhval{main/supervised-scratch/n10/test_acc/mean:3}, std \rhval{abl_supaug/supervised-+-aug/n10/test_acc/std:3}, $p=\rhval{cmp/cmp_supaug/supervised-scratch/n10/test_acc/paired_p:2}$) and remains below SimCLR at 10 labels ($p=\rhval{cmp/cmp_supaug/simclr-style-probe/n10/test_acc/paired_p:2}$; Table~\ref{tab:extra}). So in this setup augmentation of the labeled examples alone does not reproduce the SimCLR advantage.

\begin{figure}[t]\centering
\includegraphics[width=\columnwidth]{sensitivity_n50.pdf}
\caption{Sensitivity at 50 labels: temperature $\tau$ (left) and pretraining epochs (right); mean $\pm$ std over 5 seeds, including the defaults from the main runs.}\label{fig:sweep}\end{figure}

\begin{table}[t]\centering\caption{Temperature and epoch sweeps at 50 labels (5 seeds). Defaults, $\tau=0.5$ and 60 epochs, are in Table~\ref{tab:main}.}\label{tab:sweep}
\resizebox{0.85\columnwidth}{!}{\begin{tabular}{lcc}
\toprule
Method & test\_acc $\uparrow$ & $n$ \\
\midrule
SimCLR tau=0.1 & \underline{0.942 $\pm$ 0.015} & 5 \\
SimCLR tau=0.2 & 0.942 $\pm$ 0.011 & 5 \\
SimCLR tau=1.0 & \textbf{0.963 $\pm$ 0.014} & 5 \\
\bottomrule
\end{tabular}
}\\[3pt]
\resizebox{0.85\columnwidth}{!}{\begin{tabular}{lcc}
\toprule
Method & test\_acc $\uparrow$ & $n$ \\
\midrule
SimCLR ep=10 & \underline{0.940 $\pm$ 0.014} & 5 \\
SimCLR ep=30 & \textbf{0.955 $\pm$ 0.011} & 5 \\
\bottomrule
\end{tabular}
}\end{table}

\textbf{Sensitivity.} A higher temperature ($\tau=1.0$: \rhval{sweep_tau/simclr-tau-1.0/n50/test_acc/mean:3}) is slightly better than the default (\rhval{main/simclr-style-probe/n50/test_acc/mean:3}) and lower values (\rhval{sweep_tau/simclr-tau-0.1/n50/test_acc/mean:3} and \rhval{sweep_tau/simclr-tau-0.2/n50/test_acc/mean:3}) are worse (Table~\ref{tab:sweep}, Fig.~\ref{fig:sweep}); the differences are about one seed standard deviation, so we make no claim about the optimum; the $\tau=1.0$ comparison is post hoc on the test split, and the default was not tuned on it. Shortening pretraining to 10 epochs gives \rhval{sweep_epochs/simclr-ep-10/n50/test_acc/mean:3} against \rhval{main/simclr-style-probe/n50/test_acc/mean:3} at 60; 30 epochs matches the default (\rhval{sweep_epochs/simclr-ep-30/n50/test_acc/mean:3}), so accuracy saturates early at this scale.
